% !TeX root = main.tex
\lecture{9}{Mon 16 Feb 2026 11:00}{Electrostatics IX: Capacitance}

\section{Earthing}
Consider two spheres, one of radius $a$ with no charge and a smaller sphere with radius $b$ and charge $q$.

If we take the two spheres and place them in contact with each other, they will share a potential as some charge from the charged sphere is transferred to the previously uncharged sphere.

This implies that (dropping the $1/4 \pi \epsilon_0$ from both sides):
\[
 \frac{q_a}{a} = \frac{q_b}{b}
\]

And the total charge must be conserved, hence:
\[
    q = q_a + q_b
\]
\[
    \frac{q_a}{a} = \frac{q_b}{b}= \frac{q - q_a}{b} \implies (a+b)q_a = aq
\]

Hence:
\[
    q_a = \frac{aq}{a+b} \approx q \text{ for: } a \gg b
\]

So, under the same $a \gg b$ condition, $q_b \approx 0$. Effectively, touching a smaller charged sphere to a (much) bigger uncharged sphere distributes the whole charge so finely across the whole system that it is effectively zero.

Formally, earthed isolated bodies share their charge with the earth, and in doing so effectively lose that change. The bodies and earth acquire a common potential, called the zero of potential. Even though the earth gains some charge, the potential of the earth does not change because $a$ is so large.

\section{Capacitance}
Playing a charge $Q$ on a conductor alters the conductor's potential by $V \propto Q$. We define the capacitance $C$ of the conductor by the equation:
\[
    \boxed{C = \frac{Q}{V}}
\]

It has units of Coulombs per Volt, aka the Farad. Anything that can hold a charge has an associated capacitance, not just a capacitor.

For example, an isolated conducting sphere of radius $R$ with charge $Q$ has:
\[
    V = \frac{Q}{4 \pi \epsilon_0 R} \implies C = \frac{Q}{V} = 4 \pi \epsilon_0 R
\]
Capacitance is an inherent property of the material/physical setup, so shouldn't have $Q$ or $V$ in its expression.

In order to get a capacitance of a single farad, you'd need a sphere of $R = 9 \times 10^9$m, this is more than $1,000$x the that of the earth. Common lab capacitors use picofarads to microfarads.


\subsection{Capacitors}
A capacitor is a system designed to store charge (therefore energy) by separately storing positive and negative charges.

Suppose some negligible charge $\delta q$ between the two conductors, whilst they are at a potential difference $V_q$. The work done here is:
\[
    \delta W = \delta q V_q
\]

We cannot simply integrate this however, as treating $V_q$ as a constant assumes that transferred quantity of charge is negligible and does not impact the overall potential difference. This is not true for non-infinitesimal charges. We cannot therefore say $W = \int_{0}^{Q} V_q \, dq$

We instead rearrange:
\[
    q = C V_q \implies dq = C dV_q
\]
And substitute for $dq$:
\[
    W = C \int_{0}^{V} V_q \, dq = \frac{1}{2}CV^2
\]
This gives the total energy stored in a capacitor:
\[
    U = \frac{1}{2}CV^2 \qquad \left(= \frac{1}{2}QV = \frac{1}{2}\frac{Q^2}{C}\right)
\]

\section{Calculating Capacitance}
The general procedure for determining the capacitance of a system is:
\begin{enumerate}
    \item Use Gauss' Law to determine the electric field as a function of Q.
    \item Integrate to get potential difference.
    \item Use $C = Q/V$.
    \item Done!
\end{enumerate}

\subsection{Parallel Plate Capacitor}
We ignore edge effects here and assume we can treat them like infinite plates with a perfectly perpendicular and uniform electric field. This approximation holds where the side lengths of the capacitor are sufficiently larger than the distance between them. 

We have two parallel plates, each of area $A$, with separation $d$. By Gauss' law:
\[
    E = \frac{\sigma}{\epsilon_0} = \frac{Q}{\epsilon_0 A}
\]

As this is constant between the plates, the potential difference is simple to calculate:
\[
    V = - \int_{d}^{0} \pmb{E} \cdot \, d \pmb{x} = Ed = \frac{qE}{\epsilon_0 A}
\]
We integrate against the field to ensure a positively signed answer, but given we only care about a magnitude we could simply do the opposite and drop the negative sign (noting why in an exam).

Therefore:
\[
    \boxed{C = \frac{Q}{V} = \frac{\epsilon_0 A}{d}}
\]

\subsection{Long Cylindrical Capacitor}
An example is a co-axial cable, with an inner and an outer conductor, separated radially with a dielectric (next lecture) between them:
\begin{figure}[H]
    \centering
    \includegraphics[width=0.75\textwidth]{figures/lec08-05.png}
     \caption{}
\end{figure}

Here, we consider some length of the cable $L$. The inner core has charge per unit length $+\lambda$ and the outer core has $- \lambda$. We want to think about the capacitance per unit length, as in a data transmission cable this is related to the data loss.

We assume the cable is infinite to disregard length effects. The inner core has radius $r_a$ and the outer conductor radius $r_b$.

We use a cylindrical Gaussian surface of length $l$ and radius $R$. The RHS of Gauss's law gives:
\[
    \frac{Q_\text{enc}}{\epsilon_0} = \frac{\lambda l}{\epsilon_0}
\]

And the electric field is constant and radial ($\perp$ to the Gauss surface) so the LHS is:
\[
    \int \pmb{E} \cdot d \pmb{s} = E 2 \pi l
\]

Hence:
\[
    E = \frac{\lambda}{\epsilon_0 2 \pi r}
\]
\[
    \pmb{E} = \frac{\lambda}{\epsilon_0 2 \pi r} \pmb{\hat{r}}
\]

Therefore the potential difference is:
\[
    V = - \int_{r_b}^{r_a} \pmb{E} \cdot \, d \pmb{r} = - \frac{\lambda}{2 \pi \epsilon_0} \left[\ln r\right]_{r_b}^{r_a}
\]

And swapping the bounds to clear the negative and get magnitude only:
\[
    V = \frac{\lambda}{2 \pi \epsilon_0} \ln \left(\frac{r_b}{r_a}\right)
\]

Hence capacitance per unit length is:
\[
    C = \frac{Q}{V} = \frac{\lambda}{V} = \frac{2 \pi \epsilon_0}{\ln(r_b / r_a)}
\]


\subsection{Concentric Spheres}
A spherical capacitor has two spherical conductors, an inner sphere with radius $r_a$ and surface charge density $+q$ and an outer sphere with $-q$ and $r_b$. We place a spherical Gaussian surface concentric to these spheres, and between them.

Applying Gauss' law (as previously done for a sphere) we get:
\[
    \pmb{E} = \frac{Q}{4 \pi \epsilon_0 r^2} \hat{r}
\]

Hence potential difference is:
\[
    V = - \int_{r_b}^{r_a} \pmb{E} \cdot \, d \pmb{r}
\]
\[
    = - \frac{Q}{4 \pi \epsilon_0} \int_{r_b}^{r_a} \frac{1}{r^2} \, dr = \frac{Q}{4 \pi \epsilon_0} \left(\frac{1}{r_a} - \frac{1}{r_b}\right)
\]
\[
    = \frac{Q}{4 \pi \epsilon_0} \frac{(r_b - r_a)}{r_a r_b}
\]
\[
    \implies \boxed{C = \frac{Q}{V} = \frac{4 \pi \epsilon_0 r_a r_b}{r_b - r_a}}
\]

Note that the single sphere case defined earlier is effectively just a generalisation of this case where the outer sphere goes off to infinity, i.e. $r_b \to \infty$.



